\chapter{Bilangan sampai 1000}\label{ch:g2:numbers}

Kelas 1 mengikat satuan menjadi puluhan (\cref{ch:g1:tens}); sekarang
sepuluh puluhan terikat menjadi satu \emph{ratusan}, dan bilangan
memanjang sampai $1000$. Tiga angka, tiga tugas: ratusan, puluhan,
satuan.

\section{Ratusan, puluhan, satuan}

\begin{definition}[Tiga tempat]\label{def:g2:numbers:places}
Satu \emph{ratusan} adalah sepuluh puluhan, yaitu sepuluh ikat berisi
sepuluh. Bilangan $346$ punya tiga bagian:
\[
346 = 3 \text{ ratusan} + 4 \text{ puluhan} + 6 \text{ satuan}
= 300 + 40 + 6 .
\]
\end{definition}

\begin{omfigure}
\begin{tikzpicture}[scale=0.85]
  \draw (0,0) grid[xstep=2.3, ystep=0.85] (6.9,1.7);
  \node[font=\small] at (1.15,1.275) {ratusan};
  \node[font=\small] at (3.45,1.275) {puluhan};
  \node[font=\small] at (5.75,1.275) {satuan};
  \node[font=\large, omDef] at (1.15,0.425) {$3$};
  \node[font=\large, omDef] at (3.45,0.425) {$4$};
  \node[font=\large, omDef] at (5.75,0.425) {$6$};
\end{tikzpicture}

{\small Tabel nilai tempat dari $346$, dibaca ``tiga ratus empat
puluh enam''. (Bilangan yang lebih besar menunggu
\cref{ch:g3:numbers}.)}
\end{omfigure}

\begin{example}[Nol menjaga tempat]\label{ex:g2:numbers:zero}
``Tiga ratus lima'' adalah $305$: tiga ratusan, \emph{tanpa} puluhan,
lima satuan. Angka $0$ menahan $5$ tetap di tempat satuan. Tanpa \omterm{def:g1:counting:zero}{nol}
itu, $35$ menjadi bilangan yang sama sekali berbeda.
\end{example}

\begin{example}[Menukar]\label{ex:g2:numbers:exchange}
Sepuluh satuan menjadi satu puluhan, dan sepuluh puluhan menjadi
satu ratusan. Jadi $10$ puluhan $= 100$, dan $2$ ratusan $= 20$ puluhan.
Bilangan $250$ memuat $25$ puluhan, bukan hanya angka $5$.
\end{example}

\section{Membandingkan dan mengurutkan}

\begin{method}[Membandingkan bilangan tiga angka]\label{met:g2:numbers:compare}
\begin{enumerate}
  \item Angka yang lebih banyak menang: $103 > 98$;
  \item panjangnya sama: bandingkan dari kiri --- ratusan dulu, lalu
        puluhan, lalu satuan. Perbedaan pertama yang menentukan:
        $472 > 468$ karena pada puluhan, $7 > 6$.
\end{enumerate}
\end{method}

\begin{example}\label{ex:g2:numbers:order}
Urutkan $530$, $305$, $503$: ratusan dulu ($5$, $3$, $5$), jadi
$305$ di depan; antara $530$ dan $503$, puluhanlah yang menentukan ($3 > 0$):
\[
305 < 503 < 530 .
\]
\end{example}

\section{Lompatan di garis bilangan}

\begin{method}[Menambah puluhan dan ratusan]\label{met:g2:numbers:jumps}
Di garis bilangan:
\begin{enumerate}
  \item $+1$ / $-1$: mengubah satuan;
  \item $+10$ / $-10$: mengubah puluhan ($362 + 10 = 372$);
  \item $+100$ / $-100$: mengubah ratusan ($362 + 100 = 462$).
\end{enumerate}
Setiap kali hanya satu angka yang berpindah --- kecuali kalau sebuah
tempat ``penuh'' lalu menyimpan ke tempat berikutnya, seperti
$195 + 10 = 205$.
\end{method}

\begin{omfigure}
\begin{tikzpicture}[scale=1.0]
  \draw[-Stealth] (-0.3,0) -- (10.6,0);
  \foreach \x/\l in {0.5/{362}, 5.5/{372}, 9.5/{472}}
    { \draw (\x,-0.1) -- (\x,0.1);
      \node[below, font=\small] at (\x,-0.12) {$\l$}; }
  \draw[-Stealth, omDef, thick] (0.5,0.28) .. controls (2.5,1.0) ..
    (5.5,0.28);
  \node[omDef, font=\small] at (3.0,1.05) {$+10$};
  \draw[-Stealth, omThm, thick] (5.5,0.28) .. controls (7.5,1.05) ..
    (9.5,0.28);
  \node[omThm, font=\small] at (7.5,1.1) {$+100$};
\end{tikzpicture}

{\small Lompatan dari $362$: lompatan kecil $+10$ mengubah puluhan,
lompatan besar $+100$ mengubah ratusan.}
\end{omfigure}

\section{Latihan}

\begin{exercise}[$\star$]\label{exo:g2:numbers:1}
Berapa ratusan, puluhan dan satuan pada: $258$; \ $470$; \ $600$; \
$909$?
\end{exercise}

\begin{exercise}[$\star$]\label{exo:g2:numbers:2}
Tulislah bilangannya: $4$ ratusan, $0$ puluhan, $7$ satuan; \ $6$
ratusan dan $3$ puluhan; \ sepuluh ratusan.
\end{exercise}

\begin{exercise}[$\star$]\label{exo:g2:numbers:3}
Tulislah dengan angka: dua ratus enam belas; \ lima ratus empat; \
tujuh ratus sembilan puluh.
\end{exercise}

\begin{exercise}[$\star$]\label{exo:g2:numbers:4}
Uraikan seperti pada \cref{def:g2:numbers:places}: $438$; \ $702$; \
$560$.
\end{exercise}

\begin{exercise}[$\star$]\label{exo:g2:numbers:5}
Salin dan lengkapi dengan $<$ atau $>$:
\[
327 \;?\; 372, \qquad
580 \;?\; 508, \qquad
199 \;?\; 200, \qquad
640 \;?\; 604 .
\]
\end{exercise}

\begin{exercise}[$\star$]\label{exo:g2:numbers:6}
Urutkan dari yang terkecil ke yang terbesar: $430$; \ $304$; \ $340$; \ $403$.
\end{exercise}

\begin{exercise}[$\star$]\label{exo:g2:numbers:7}
Hitunglah dengan melompat di garis bilangan: $253 + 10$; \quad
$253 + 100$; \quad $480 - 10$; \quad $716 - 100$.
\end{exercise}

\begin{exercise}[$\star$]\label{exo:g2:numbers:8}
Lanjutkan membilang: $195, 205, 215, ?, ?$. \quad
$300, 400, 500, ?, ?$. \quad Hati-hati dengan yang pertama!
\end{exercise}

\begin{exercise}[$\star$]\label{exo:g2:numbers:9}
Sebuah toko punya $6$ kotak berisi $100$ pensil dan $5$ pensil lepas.
Berapa pensil seluruhnya?
\end{exercise}

\begin{exercise}[$\star\star$]\label{exo:g2:numbers:10}
Dengan memakai angka $2$, $5$ dan $8$ masing-masing sekali, tulislah
bilangan tiga angka yang terbesar dan yang terkecil. Berapa jarak
keduanya? (Kamu boleh menghitungnya di garis bilangan dengan lompatan.)

\end{exercise}
